% !PRINTED 731
\weberpage{731}

From this it can be concluded, as in \S.\,187, that there are systems of $\nu'$ independent normal units whose regulator $T_0$ becomes \gesperrt{as small as possible}, and such systems are called \gesperrt{fundamental systems of normal units}.

If $\mathcal{E}_0, \mathcal{E}_1, \ldots, \mathcal{E}_{\nu'-1}$ is such a fundamental system, then the exponents $\xi_0, \xi_1, \ldots, \xi_{\nu'-1}$ in (10) and (11) come out as \gesperrt{integral rational numbers}. This too can be proved in the same way as in \S.\,187.

If, under this assumption about the $\mathcal{E}_i$, we set up the equations (11) for any other system of independent units $\mathcal{E}_0^{(1)}, \mathcal{E}_1^{(1)}, \ldots, \mathcal{E}_{\nu'-1}^{(1)}$, then we obtain expressions of the form:

\[
L_{i,k}^{(1)} = \xi_{i,0} L_{0,k} + \xi_{i,1} L_{1,k} + \cdots + \xi_{i,\nu'-1} L_{\nu'-1,k}.
\]

If, then, $T_1$ is the regulator of the system of the $\mathcal{E}_i^{(1)}$, it follows by the multiplication theorem for determinants, when by $C$ is denoted the absolute value of the determinant of the integer exponents $\xi_{i,k}$, and hence an \gesperrt{integral rational positive number}:

\weq{13}{T_1 = C\,T_0.}

It cannot be proved, and for the higher values of $m$ it probably does not even hold, that $\tau_0, \tau_1, \ldots, \tau_{\nu'-1}$ form a fundamental system\footnote{It is here, as with all questions that concern the units, very difficult to arrive at definite numerical results. According to the results contained in the comprehensive work of \gesperrt{Reuschle}, \glqq Tafeln complexer Primzahlen, welche aus Wurzeln der Einheiten gebildet sind\grqq{} (Berlin 1875), the $\tau$ still form a fundamental system for $m=16$. For here the fundamental number of the field $H_m$ is equal to $2^{11}$, the degree of the field is equal to 4, and hence, by the remark to \S.\,156, there is in every ideal class an ideal whose norm is less than 6. But according to the tables of \gesperrt{Reuschle} every natural prime below 100 can be decomposed in the field $H_m$ into actually existing prime factors. Consequently only principal ideals certainly divide all numbers below 6, and consequently there is here only the one class, the principal class. Since, as we have seen, the first class-number factor is also $=1$, the cyclotomic field $\Omega_{16}$ is of one class. In it the same laws of primes hold as in the field of rational numbers. This is doubtful for the field $\Omega_{32}$, and certainly no longer true in the field $\Omega_{64}$, in which the class number is a multiple of 17.}.

For our purpose, however, it is of importance that these units, with respect to the denominator 2 in the exponents
%CONT>>
